\documentclass[11pt,a4paper]{article}
\usepackage[margin=25mm]{geometry}
\usepackage[T1]{fontenc}
\usepackage[utf8]{inputenc}
\usepackage{lmodern}
\usepackage{amsmath,amssymb,amsthm}
\usepackage{booktabs,array,longtable}
\usepackage{microtype}
\usepackage{fancyhdr}
\usepackage{xurl}
\usepackage[hidelinks]{hyperref}
\hypersetup{pdftitle={Congruence carriers for Diophantine consistency equations},
 pdfauthor={Research prepared for Vladimir Reshetnikov and ProveIt}}
\setlength{\headheight}{14pt}
\pagestyle{fancy}
\fancyhf{}
\fancyhead[L]{\small Congruence carriers and exact eliminations}
\fancyhead[R]{\small ProveIt research}
\fancyfoot[C]{\thepage}
\setlength{\parskip}{4pt}
\setlength{\parindent}{0pt}
\setlength{\emergencystretch}{2em}
\setcounter{tocdepth}{1}
\allowdisplaybreaks[2]
\newtheorem{theorem}{Theorem}[section]
\newtheorem{lemma}[theorem]{Lemma}
\newtheorem{proposition}[theorem]{Proposition}
\newtheorem{corollary}[theorem]{Corollary}
\theoremstyle{definition}
\newtheorem{definition}[theorem]{Definition}
\theoremstyle{remark}
\newtheorem{remark}[theorem]{Remark}
\newcommand{\Z}{\mathbb Z}
\newcommand{\Zpos}{\mathbb Z_{>0}}
\newcommand{\Con}{\operatorname{Con}}
\newcommand{\ZFC}{\mathsf{ZFC}}
\newcommand{\pos}{\operatorname{Pos}}
\newcommand{\ini}{\operatorname{in}}
\newcommand{\chiA}{\chi_A}
\newcommand{\psiA}{\psi_A}
\newcommand{\sourcepath}[1]{\texttt{\small #1}}
\newcommand{\snapshot}{\texttt{beb0192cb8dc36f78c369c080f2adfbc1887a486}}
\title{\textbf{Congruence carriers for Diophantine consistency equations}\\[5pt]
 \large Thirteen positive unknowns in degree 34,\\
 twelve positive unknowns in degree 118,\\
 and two reductions over the integers}
\author{Research prepared for Vladimir Reshetnikov and the ProveIt repository\\
 \small with OpenAI assistance}
\date{23 October 2026}
\begin{document}
\maketitle
\begin{abstract}
We give two fixed integer polynomials whose lack of strictly positive
integer zeros is provably equivalent in ZFC to the syntactic consistency
of ZFC. The first has thirteen unknowns, exact total degree \(34\), and
a literal evaluator with \(93\) binary arithmetic operations. The second
has twelve unknowns, exact degree \(118\), and \(104\) operations.
Their coefficients are the complete, explicitly materialized compact
native-compiler tuple already present in ProveIt.

The main construction puts divisibility information into two existing
Pell square equations. The first equation simultaneously recovers a
scale and a transport quotient, with positivity established before any
Pell-index argument. The second carries two congruences in a prescribed
order: an index is recovered modulo \(q^4\), then a missing scale is
recovered modulo \(q^5\). A divided auxiliary norm removes one more
positive unknown. We supply the signed index lemma, all new positivity
arguments, and explicit canonical completeness constructions.

Separately, exact substitutions in a fixed forty-five-integer quartic
source give equations in \(37\) unrestricted integer variables of degree
\(6\), and \(34\) unrestricted integer variables of degree \(8\). Those
two equations retain their source's different complete coefficient tuple.
The archive includes literal equations, integer coefficient payloads,
reconstruction programs, and exact algebraic checks. The proof composes
ordinary mathematical reductions with named inherited compiler and
refutation-calibration theorems. It is not a new end-to-end
proof-assistant certificate, a consistency proof, or an optimality theorem.
\end{abstract}
\thispagestyle{empty}
\clearpage
\tableofcontents
\clearpage

\section{Results, domains, and the comparison being made}
\label{sec:results}

Fix the recursive presentation of ZFC and the finite-proof coding used by
the inherited ProveIt calibration. Write \(\Con(\ZFC)\) for the arithmetic
assertion that this proof system has no finite syntactic refutation.
The objective is a fixed polynomial with integer coefficients; there is
no variable exponent, parameter selected by a consistency assumption,
or uncounted existential register in any polynomial below.

\begin{theorem}[Main constructions]
\label{thm:main}
There are explicitly specified polynomials \(P_{13}\) and \(P_{12}\), with
the complete coefficient tuple of Section~\ref{sec:binding}, such that
\begin{align*}
 \ZFC&\vdash
 \bigl(\neg\exists\boldsymbol{x}\in\Zpos^{13}\ 
                  P_{13}(\boldsymbol{x})=0\bigr)
       \ \longleftrightarrow\ \Con(\ZFC),\\
 \ZFC&\vdash
 \bigl(\neg\exists\boldsymbol{x}\in\Zpos^{12}\ 
                  P_{12}(\boldsymbol{x})=0\bigr)
       \ \longleftrightarrow\ \Con(\ZFC).
\end{align*}
The exact degree and literal arithmetic-operation pairs are \(34/93\)
and \(118/104\), respectively. There are also explicitly specified
polynomials \(Q_{37}\) and \(Q_{34}\) over unrestricted integers with
the same consistency equivalence, in \(37\) variables of degree \(6\)
and \(34\) variables of degree \(8\).
\end{theorem}

The meaning of the displayed provability claim is ordinary mathematical
provability formalizable in ZFC, using the inherited compiler and
calibration theorem stated in Section~\ref{sec:source}. The archive does
not export a newly checked object-language ZFC derivation or a complete
Lean/Rocq build of that composition. The new polynomial identities,
divisibility arguments, and existence reductions are proved here; the
remaining source boundary is explicit.

\subsection{A fixed, bounded comparison}

All comparisons refer to ProveIt revision \snapshot.
The relevant source directories are 44, 57, 58, and 59 of
\path{Computability/DiophantineConsistency/con-zfc-diophantine-forms}.
The principal improvements and their costs are as follows.

\begin{center}
\begin{tabular}{@{}lrrrr@{}}
\toprule
Construction & Domain & Variables & Degree & Operations\\
\midrule
Source 59 & \(\Zpos\) & 13 & 126 & 91\\
\(P_{13}\), this paper & \(\Zpos\) & 13 & 34 & 93\\
Source 57 & \(\Zpos\) & 12 & 242 & 94\\
\(P_{12}\), this paper & \(\Zpos\) & 12 & 118 & 104\\
\midrule
Historical sextic & \(\Z\) & 42 & 6 & ---\\
\(Q_{37}\), this paper & \(\Z\) & 37 & 6 & ---\\
Historical octic & \(\Z\) & 36 & 8 & ---\\
\(Q_{34}\), this paper & \(\Z\) & 34 & 8 & ---\\
\bottomrule
\end{tabular}
\end{center}

The positive-variable comparison is to the explicit source-57 and
source-59 evaluators \cite{source57,source59}. A smaller degree need not
give a smaller arithmetic circuit. The earlier thirteen-positive point
of degree \(132\) at \(83\) operations, and twelve-positive point of
degree \(258\) at \(89\) operations, retain lower operation counts
\cite{source57}. The integer comparisons are to the preserved historical
degree-six and degree-eight packets discussed in
Section~\ref{sec:integer}; the octic here uses two more residual equations
than its historical comparator.

These are separate tradeoffs. In particular, a polynomial over
\(\Zpos\) is not being advertised as a polynomial with the same number
of unrestricted integer unknowns. Nor is a degree/variable comparison an
operation-count comparison. Existing small-variable equations at much
larger degrees remain different points in the catalogue.

\subsection{Size conventions}

Total degree counts the supplied unknowns and assigns degree zero to the
four fixed coefficient ports. A named register denotes its fully
substituted integer polynomial, not a new variable. The operation counts
refer to the supplied straight-line graphs: every live binary
addition, subtraction, and multiplication is charged, including
multiplication by a small numeral. They are reproducible costs of those
evaluators, not lower bounds on arithmetic complexity.

The large fixed integers are delivered as byte-exact payloads and by
deterministic reconstruction programs. Their bit lengths are reported
separately. A compact circuit does not make those coefficients small, and
an expanded polynomial is not needed to specify them unambiguously.

\section{The inherited source interface}
\label{sec:source}

The main positive-variable constructions use one whole authentic tuple
\[
 (\beta,K_0,M_C,M_F),\qquad
 B=2^{590069},\qquad \beta=B-1.
\]
Here \(M_F\) is the already shifted high-mask port:
\(M_F=M_F^{\rm native}+B-1\). It must not be shifted a second time.
The ports are not independently adjustable. The source-44 adapter proves
its native interface for their joint values \cite{source44}.

\subsection{The normalized positive kernel}

The old positive coordinates are
\[
 J,F,\alpha,z,f,h,i,T_{\rm old},s,w,\tau,\eta,\zeta,
 y,\gamma_{\rm old}.
\]
For this subsection only, define
\begin{align}
 q&=\beta J+1,& C&=q-F-\alpha,&
 X&=wq,& Y&=sq^3, \nonumber\\
 E&=XY,& k&=\eta+\zeta,& c&=kY+\eta,&
 a&=E+Y, \nonumber\\
 H&=4a+3,& \Delta&=a^2+H=(a+2)^2-1,&
 \ell&=EkY,& I&=ic^2.                     \label{eq:oldreg}
\end{align}
Further put
\begin{align}
 R&=\bigl((q-1)(q-F)+\alpha\bigr)(q^2-1)
                     +(M_C+qM_F)J, \label{eq:R}\\
 D&=X+ac+H\gamma_{\rm old}, \nonumber\\
 V_{\rm old}&=c(T_{\rm old}f-1)-Rf^2.       \label{eq:oldV}
\end{align}
The normalized positive source system is
\begin{align}
 \tau^2-\ell(\ell+k)&=1, &
 D^2-\Delta c^2&=1, \label{eq:oldnorms}\\
 f^2-\Delta I^2&=1, &
 \Delta(f^2-1)(V_{\rm old}^2-y^2)+y^2&=1,
                                                        \label{eq:oldaux}\\
 k-hE-R&=1, &
 (K_0+w)C+(q-F)-z(q-1)&=1.                \label{eq:oldtransport}
\end{align}

\begin{theorem}[Inherited compiler/calibration interface]
\label{thm:source}
For the whole compact source-44 tuple, a strictly positive solution of
\eqref{eq:oldreg}--\eqref{eq:oldtransport} yields a finite syntactic
refutation of ZFC. Conversely, a finite refutation yields canonical
strictly positive solutions of this system with the properties stated
below.
\end{theorem}

This theorem is the composition of the native-window/Pell compiler,
its compact-coefficient adapter, and the literal blank-K215
halting/refutation calibration \cite{source44,ordinary,pellaudit,source58}.
The normalized strong coordinate used here is converted to the ordinary
strong coordinate by \(i_{\rm ordinary}=\Delta i\). Then
\[
 (i_{\rm ordinary}c^2)^2=\Delta(f^2-1),
\]
exactly matching the ordinary kernel. Thus the stated interface does not
assume that a theorem for an arbitrary coefficient tuple applies to a
different numerical instance.

\subsection{Outer bounds available before index recovery}
\label{sec:outer}

The authentic masks satisfy
\[
 0<M_C,M_F^{\rm native}<B-1,\qquad K_0>0.
\]
Consequently, if \(J,F,\alpha>0\) and \(C=q-F-\alpha>0\), then
\begin{equation}
 q\ge16,\qquad 3q<R<q^4-2.                 \label{eq:outerbounds}
\end{equation}
These inequalities use no divisibility assumption on directly supplied
\(X,Y\). Indeed, with \(G=(q-1)(q-F)+\alpha\),
\[
 G=q\alpha+(q-1)C,\qquad q\le G\le q^2-q,
\]
and
\[
 0<(M_C+qM_F)J<(q-1)(1+2q).
\]
The resulting upper bound is
\(R<q^4-q^3+q^2-1<q^4-2\); the lower bound is
\(R>q(q^2-1)>3q\). The distinction matters in our construction:
\(C>0\) is proved before \eqref{eq:outerbounds} is used, whereas
\(q\mid X\) is recovered substantially later.

\subsection{Canonical promises used in completeness}
\label{sec:canonical}

The source-44 padding controller and source-58 interface provide canonical
witnesses with
\begin{align}
 q&=B^N,& J&=1+B+\cdots+B^{N-1},&
 N&=581197^e>1,                              \label{eq:repunit}\\
 R&\equiv3\pmod4,& X&=2^R,&
 c&=\psi_{a+2}(R),                            \label{eq:canonicalR}\\
 \tau&=\chi_{P_0}(n_0),&
 k&=2\psi_{P_0}(n_0),&
 P_0&=2XY^2+1,\quad n_0=(R+1)/2\ge3.          \label{eq:canonicalpell}
\end{align}
The Pell notation is defined in Section~\ref{sec:pell}.
In \eqref{eq:repunit}, \(N\) is odd, and \(e\) is a proof-only padding
exponent. These are canonical completeness promises, not restrictions
silently imposed on arbitrary solutions.

The compact adapter proves the controller at the actual final packed
index, after its permitted dummy choices, and retains the low-mask
identity \(R\equiv3\pmod4\) \cite[Section 6]{source44}.
We will also use
\begin{equation}
 \gamma_{\rm old}>X\ge w.                    \label{eq:canonicalgamma}
\end{equation}
Here is the source-58 growth proof in the present notation.
Since \(D=\sqrt{\Delta c^2+1}>(a+1)c\),
\[
 \gamma_{\rm old}=\frac{D-ac-X}{H}>\frac{c-X}{H}.
\]
From \(n_0\ge3\),
\[
 k\ge2\psi_{P_0}(3)=8P_0^2-2>32X^2Y^4,\qquad
 c>kY>32X^2Y^5.
\]
As \(X\ge16\), \(Y\ge1\), and \(H=4XY+4Y+3\le9XY\), these inequalities
give \(c-X>XH\), proving \eqref{eq:canonicalgamma}. Merely knowing
\(\gamma_{\rm old}>0\) would not justify the subtraction of \(w\) later.

\section{Pell lemmas, including the signed auxiliary index}
\label{sec:pell}

For an integer \(A\ge2\), put \(\Delta=A^2-1\) and define
\[
 \chi_A(0)=1,\quad \chi_A(1)=A,\qquad
 \psi_A(0)=0,\quad \psi_A(1)=1,
\]
both sequences satisfying \(z_{n+2}=2Az_{n+1}-z_n\).
Equivalently,
\[
 \chi_A(n)+\psi_A(n)\sqrt{\Delta}
       =(A+\sqrt{\Delta})^n.
\]
We use their addition identities, positive-solution classification, and
strong divisibility
\[
 \gcd(\psi_A(r),\psi_A(s))=\psi_A(\gcd(r,s)).
\]
These standard Pell facts and their use by the source are collected in
the inherited audit \cite{pellaudit}. The positive fundamental
integer-coordinate solution for discriminant \(A^2-1\) is \((A,1)\).

\begin{lemma}[Rational square-root integrality]
\label{lem:rational}
If a rational number has integer square, it is an integer. In particular,
if \(d>0\) and \(z^2=d^2b\), with \(z,b,d\in\Z\), then \(d\mid z\).
\end{lemma}
\begin{proof}
Write the rational number as \(a/b\) in lowest terms with \(b>0\).
If \(a^2/b^2\) is an integer, then \(b^2\mid a^2\), so coprimality
forces \(b=1\).
\end{proof}

\begin{lemma}[Signed abscissa step-down]
\label{lem:stepdown}
Let \(A\ge2\), \(m\ge1\), and \(0<k<m\). If \(t\ge0\) and
\[
 \chi_A(t)\equiv\chi_A(k)\pmod{\chi_A(m)},
\]
then \(t\equiv k\) or \(-k\pmod{2m}\).
\end{lemma}
\begin{proof}
Write \(F_m=\chi_A(m)\). The addition formulas give
\[
 \chi_A(j+2m)\equiv-\chi_A(j),\qquad
 \chi_A(2m-j)\equiv-\chi_A(j)\pmod{F_m}.
\]
Indeed \(\chi_A(2m)=2F_m^2-1\) and
\(\psi_A(2m)=2F_m\psi_A(m)\). Reducing \(t\) modulo \(2m\), and
reflecting a representative above \(m\), therefore gives a signed
representative \(\varepsilon\chi_A(s)\), \(0\le s\le m\).
The case \(s=m\) has residue zero, while
\(0<\chi_A(k)<F_m\). If \(s<m\) and \(\varepsilon=-1\), the
congruence would force
\[
 F_m\mid\chi_A(s)+\chi_A(k),\qquad
 0<\chi_A(s)+\chi_A(k)\le2\chi_A(m-1)<F_m,
\]
a contradiction. The last inequality follows from the recurrence for
\(A\ge2\). For the positive sign, both representatives are strictly
between zero and \(F_m\); their equality modulo \(F_m\) is equality as
integers. Strict growth gives \(s=k\). The index operations used only
sign and multiples of \(2m\), proving the conclusion.
\end{proof}

This is the explicit signed argument of the ordinary source kernel
\cite[Section 5]{ordinary}, also used by its normalized descendant
\cite[Section 5]{source38}. The hypothesis is \(k<m\), not an unproved
smaller half-interval for \(k\).

\begin{lemma}[Odd-index quotient identities]
\label{lem:odd}
There are integer polynomials \(Q_v\) such that
\begin{align}
 \chi_Z(2v+1)&=ZQ_v(Z^2), \nonumber\\
 Q_v(0)&=(-1)^v(2v+1),&
 Q_v(1-A^2)&=(-1)^v\psi_A(2v+1).             \label{eq:oddidentities}
\end{align}
Moreover
\begin{equation}
 \psi_Z(2v+1)\equiv(-1)^v\pmod{Z^2}.        \label{eq:oddordinate}
\end{equation}
\end{lemma}
\begin{proof}
For the quotient polynomials,
\[
 Q_0(x)=1,\quad Q_1(x)=4x-3,\quad
 Q_{v+1}(x)=(4x-2)Q_v(x)-Q_{v-1}(x).
\]
The odd subsequence of the Pell recurrence proves the first identity.
Evaluation at zero proves the first evaluation by induction.
At \(x=1-A^2\), the recurrence coefficient changes sign relative to
the odd \(\psi_A\) subsequence, whose initial values are
\(1,4A^2-1\); this proves the second evaluation.
For \eqref{eq:oddordinate}, write \(b_v=\psi_Z(2v+1)\).
Then \(b_0=1\), \(b_1=4Z^2-1\), and
\(b_{v+1}=(4Z^2-2)b_v-b_{v-1}\). Modulo \(Z^2\), induction gives
\(b_v=(-1)^v\).
\end{proof}

\begin{lemma}[Signed auxiliary-index recovery]
\label{lem:index}
Let \(A\ge2\), \(\Delta=A^2-1\), and let positive integers
\(p,m,R,c,f,S\) satisfy
\[
 c=\psi_A(p),\quad f=\chi_A(m),\quad
 c\mid m,\quad m>2p,\quad S\ge2,\quad c\mid S,\quad
 S^2=\Delta(f^2-1),
\]
and \(0<p,R<c/2\). If integers \(V,y\) satisfy
\begin{equation}
 S^2V^2-(S^2-1)y^2=1,\qquad
 V\equiv-c\pmod f,\qquad V\equiv-R\pmod c,    \label{eq:indexhyp}
\end{equation}
then \(p=R\). The signs of \(V,y\) are unrestricted.
\end{lemma}
\begin{proof}
The norm equation excludes \(V=0\) and \(y=0\).
Positive Pell classification at parameter \(S\) gives \(\nu\ge1\) with
\[
 S|V|=\chi_S(\nu),\qquad |y|=\psi_S(\nu).
\]
The index \(\nu\) is odd: for even indices,
\(\chi_S(2v)\equiv(-1)^v\pmod S\), contrary to divisibility by \(S\).
Write \(\nu=2v+1\). Then \(|V|=Q_v(S^2)\).
Since
\(S^2=\Delta(f^2-1)\equiv1-A^2\pmod f\), the first congruence in
\eqref{eq:indexhyp} and Lemma~\ref{lem:odd} give
\(\psi_A(\nu)\equiv\pm c\pmod f\). Squaring and using
\(\chi_A(2j)=1+2\Delta\psi_A(j)^2\) gives
\[
 \chi_A(2\nu)\equiv\chi_A(2p)\pmod{\chi_A(m)}.
\]
Lemma~\ref{lem:stepdown} applies with \(k=2p<m\), and yields
\(\nu\equiv\pm p\pmod m\), hence modulo \(c\).

On the other hand, \(c\mid S\) and the second evaluation in
\eqref{eq:oddidentities} at zero give
\(\nu\equiv\pm R\pmod c\). Thus \(p\equiv R\) or \(-R\pmod c\).
The strict inequalities \(0<p,R<c/2\) exclude the negative
identification and every nonzero multiple of \(c\). Hence \(p=R\).
\end{proof}

No main-root congruence, deleted index quotient, or positive sign of \(V\)
is a hypothesis of this lemma. A signed strong root may be replaced by
its absolute value as the modulus, with no change to its square or its
congruences. The lemma is the exact coefficient-independent implication
in \cite[Section 5]{ordinary,source57}, now stated without compiler
notation.

\begin{lemma}[Direct normalized rank]
\label{lem:rank}
Suppose \(c=\psi_A(p)>2\) and positive integers \(f,I\) satisfy
\[
 f^2-\Delta I^2=1,\qquad c^2\mid I.
\]
Then \(f=\chi_A(m)\), \(I=\psi_A(m)\) for some \(m\ge1\), and
\[
 pc\mid m,\qquad c\mid m,\qquad m>2p.
\]
\end{lemma}
\begin{proof}
Pell classification gives \(m\). Strong divisibility gives \(p\mid m\).
Write \(m=pv\) and \(D=\chi_A(p)\). Composition gives
\(I=c\psi_D(v)\). Since \(D^2\equiv1\pmod c\), the recurrence gives
\[
 \psi_D(v)\equiv vD^{v-1}\pmod c.
\]
The number \(D\) is invertible modulo \(c\). Since \(c^2\mid I\),
we obtain \(c\mid v\); hence \(pc\mid m\), and \(m\ge pc>2p\).
\end{proof}

This argument is the normalized bridge in
\cite[Section 3.1]{pellaudit} and \cite[Section 5]{source57}.
It needs no estimate \(c>A\Delta^2\) from the older ordinary-strong rank
lemma. Avoiding that estimate is essential when the first index quotient
has not yet been restored.

\begin{lemma}[Positive auxiliary reconstruction]
\label{lem:auxcomplete}
Suppose \(R\ge3\), \(R\equiv3\pmod4\), \(c=\psi_A(R)\), and
\[
 f^2-\Delta I^2=1,\qquad f,I>0,\qquad c\mid I.
\]
Set \(S=\Delta I\), \(Q=S^2\). There are positive integers \(T,u\) with
\[
 V=f^2(cT-R)-c,\qquad
 QV^2-(Q-1)(V+u)^2=1.
\]
They may be chosen with \(y=V+u\equiv-1\pmod Q\).
\end{lemma}
\begin{proof}
Write \(R=2v+1\), where \(v\) is odd, and define
\[
 V_*=\frac{\chi_S(R)}S=Q_v(S^2),\qquad y_*=\psi_S(R).
\]
They are positive integers satisfying
\[
 QV_*^2-(Q-1)y_*^2=1.
\]
Since \(S^2\equiv1-A^2\pmod{f^2}\) and \(c\mid S\),
Lemma~\ref{lem:odd} gives
\[
 V_*\equiv-c\pmod{f^2},\qquad V_*\equiv-R\pmod c,\qquad
 y_*\equiv-1\pmod Q.
\]
The strong equation gives \(f^2\equiv1\pmod c\), so
\(\gcd(c,f^2)=1\). Therefore
\begin{equation}
 T=\frac{V_*+c+Rf^2}{cf^2}\in\Zpos.          \label{eq:Tcomplete}
\end{equation}
It has \(f^2(cT-R)-c=V_*\).
Since \(R\ge3\), \(S\ge2\), we have \(V_*\ge4S^2-3>1\), and
\[
 y_*^2-V_*^2=\frac{V_*^2-1}{Q-1}>0.
\]
Thus \(u=y_*-V_*>0\). This proves all claims.
\end{proof}

This is the squared-modulus reconstruction of source 51 and
source 58 \cite{source51,source58}. It chooses fresh auxiliary witnesses;
it does not assert that every old quotient was divisible by \(f\), or
that every arbitrary auxiliary ordinate has the displayed residue.

\section{The thirteen-positive polynomial}
\label{sec:p13}

Supply exactly the thirteen unknowns
\begin{equation}
 j,F,\alpha,f,i,T,X,Y,r,\eta,\zeta,u,\gamma\in\Zpos.
                                                        \label{eq:vars13}
\end{equation}
All subsequent letters are polynomial abbreviations. In particular,
\(X,Y\) are supplied directly and are not initially assumed to be
divisible by any power of \(q\). Define
\begin{align}
 J&=3j+1,& q&=\beta J+1,& C&=q-F-\alpha, \nonumber\\
 E&=XY,& k&=\eta+\zeta,& c&=kY+\eta, \nonumber\\
 a&=E+Y,& H&=4a+3,& \Delta&=a^2+H, \nonumber\\
 \ell&=EkY,& b_0&=X+ac,& I&=ic^2.             \label{eq:common13}
\end{align}
\begin{align}
 R&=\bigl((q-1)(q-F)+\alpha\bigr)(q^2-1)
                             +(M_C+qM_F)J, \nonumber\\
 r_0&=k-R-1,\qquad U=q^5, \nonumber\\
 D_*&=Ub_0+H(U\gamma+q^4X-r_0),              \label{eq:maincarrier}\\
 K_a&=\Delta(f^2-1),\qquad V=f^2(cT-R)-c,     \label{eq:aux13}\\
 W&=(q-1)k+1,\qquad M_0=q^3(q-1), \nonumber\\
 B_t&=(qK_0+X)C-qF, \nonumber\\
 Z&=M_0(\ell+k+r)+(q-1)Y-q^3WB_t.            \label{eq:firstcarrier}
\end{align}
The four residuals are
\begin{align}
 A_0&=Z^2-M_0^2\bigl(\ell(\ell+k)+1\bigr),     \label{eq:A0}\\
 A_1&=D_*^2-U^2(1+\Delta c^2),               \label{eq:A1}\\
 A_a&=(V+u)^2-K_a u(2V+u)-1,                 \label{eq:Aa}\\
 A_s&=f^2-\Delta I^2-1.                       \label{eq:As}
\end{align}
Our first fixed polynomial is
\begin{equation}
 \boxed{P_{13}=A_0^2+A_1^2+A_a^2+A_s^2.}     \label{eq:P13}
\end{equation}
Equations \eqref{eq:common13}--\eqref{eq:P13}, together with the coefficient
binding, are a literal division-free specification. The archive graph
\sourcepath{positive/p13\_degree34.json} gives its shared-register
evaluation.

\subsection{Why these two carriers are useful}

An equation \(z^2=m^2b\), with integer \(b\), forces the rational root
\(z/m\) to be integral. In \(A_0\), the coprime factors \(q^3\) and
\(q-1\) recover two different quotients. The sign of its recovered root
is allowed to vary; the bound on that root forces the transport numerator
to be positive, and therefore forces \(C>0\).

The main square \(A_1\) then gives
\[
 q^5\mid q^4X-(k-R-1).
\]
Its reduction modulo \(q^4\) is used first, after the auxiliary Pell index
has been recovered. Only then is divisibility by \(q^5\) used to obtain
\(q\mid X\). The two stages are necessary: \(X=wq\) is unavailable during
the preliminary index bounds. The normalized strong equation is retained
so those bounds can be proved without the removed index quotient.

\section{Soundness of the thirteen-variable equation}
\label{sec:sound13}

At a positive zero of \(P_{13}\), every residual in
\eqref{eq:A0}--\eqref{eq:As} vanishes. We reconstruct the source witnesses
in an order that establishes each denominator and positivity premise
before it is used.

\subsection{The first carrier proves positivity before index recovery}

By Lemma~\ref{lem:rational}, \(A_0=0\) gives
\[
 \tau=\frac Z{M_0}\in\Z,\qquad
 \tau^2=\ell(\ell+k)+1.
\]
The integer \(\tau\) may be negative. Modulo \(q^3\), the divisibility
\(M_0\mid Z\) gives \(q^3\mid(q-1)Y\), hence \(q^3\mid Y\).
Modulo \(q-1\), it gives \((q-1)\mid B_t\), because
\(q^3\equiv W\equiv1\pmod{q-1}\).
Define the proof-only integers
\[
 s=Y/q^3>0,\qquad t=B_t/(q-1).
\]
Dividing the carrier by \(M_0\) yields
\begin{equation}
 \tau=\ell+k+r+s-Wt.                         \label{eq:transportroot}
\end{equation}
Since
\[
 (\ell+k)^2-\bigl(\ell(\ell+k)+1\bigr)
       =\ell k+k^2-1>0,
\]
we have \(|\tau|<\ell+k\). Thus
\[
 Wt=\ell+k+r+s-\tau>r+s>0.
\]
The register \(W\) is positive, so \(t>0\) and \(B_t>0\).
If \(C\le0\), however,
\[
 B_t=(qK_0+X)C-qF<0,
\]
which is impossible. We have proved \(C>0\) without using \(q\mid X\)
or any auxiliary index theorem. The bounds
\[
 3q<R<q^4-2
\]
now follow from Section~\ref{sec:outer}.

\subsection{The second carrier gives a positive main root}

As \(590069\) is odd, \(B\equiv2\pmod3\) and \(\beta\equiv1\pmod3\).
The imposed \(J=3j+1\) gives \(q\equiv2\pmod3\).
Since \(q^3\mid Y\),
\[
 H=4Y(X+1)+3\equiv3\pmod q,\qquad \gcd(q,H)=1.
\]
The numerator \(D_*\) is positive before we know \(q\mid X\).
Indeed \(r_0<k\), \(c>kY\), and
\[
 UaY>H
\]
because \(U=q^5\), \(q\ge16\), \(Y\ge q^3\), and \(H=4a+3<6a\).
Therefore
\begin{equation}
 D_*>Uac-Hk>k(UaY-H)>0.                     \label{eq:Dpositive}
\end{equation}
From \(A_1=0\), Lemma~\ref{lem:rational} gives
\begin{equation}
 D=D_*/U\in\Zpos,\qquad D^2-\Delta c^2=1.
                                                        \label{eq:Drecover}
\end{equation}
Reducing the numerator modulo \(U=q^5\), and cancelling \(H\) using
\(\gcd(H,q)=1\), gives
\begin{equation}
 q^5\mid q^4X-r_0.                          \label{eq:q5carrier}
\end{equation}

\subsection{Preliminary Pell bounds with directly supplied scales}

Put \(A=a+2\) and \(P_0=2XY^2+1\). Main Pell classification gives
\[
 D=\chi_A(p),\qquad c=\psi_A(p),\qquad p\ge1.
\]
For the first norm, writing \(v=XY^2\) gives
\(\tau^2-v(v+1)k^2=1\).
Its positive fundamental integer-coordinate unit is \((2v+1,2)\):
ordinate \(1\) would require \(v^2+v+1\) to be a square strictly between
\(v^2\) and \((v+1)^2\), while ordinate \(2\) occurs.
Consequently
\[
 |\tau|=\chi_{P_0}(n_0),\qquad
 k=2\psi_{P_0}(n_0),\qquad n_0\ge1.
\]
Since \(c>1\), \(p\ge2\). If \(n_0=1\), then \(k=2\),
\(\eta=\zeta=1\), and \(c=2Y+1\). This contradicts
\[
 c\ge\psi_A(2)=2A\ge4Y+4.
\]
Thus \(n_0\ge2\), and
\begin{equation}
 k\ge4P_0>8XY^2,\qquad
 c>kY>8XY^3\ge8q^9>2R.                     \label{eq:cR}
\end{equation}
Here \(X\ge1\), not \(X\ge q\), is sufficient.
Also \(P_0>A\) and \(c>k\); monotonicity in parameter and index yields
\(n_0<p\). The recurrence, with \(A\ge3\), gives
\begin{equation}
 c=\psi_A(p)\ge(2A-1)^{p-1}\ge5^{p-1}>2p
 \qquad(p\ge2).
                                                        \label{eq:cp}
\end{equation}
Thus both half-interval hypotheses for Lemma~\ref{lem:index} are
available before index recovery.

\subsection{Normalized rank and the exact main index}

From \(A_s=0\), the normalized strong equation is
\[
 f^2-\Delta I^2=1,\qquad I=ic^2.
\]
Lemma~\ref{lem:rank} provides an index \(m\) with
\[
 f=\chi_A(m),\qquad I=\psi_A(m),\qquad
 c\mid m,\quad m>2p.
\]
Set \(S=\Delta I\). Then \(S^2=\Delta(f^2-1)=K_a\), and \(c\mid S\).
With \(y=|V+u|\), the auxiliary residual becomes
\begin{equation}
 S^2V^2-(S^2-1)y^2=1.                       \label{eq:auxrecover13}
\end{equation}
It excludes \(y=0\) and \(V=0\). The formula for \(V\) gives
\[
 V\equiv-c\pmod f,\qquad
 V\equiv-R\pmod c,
\]
the latter because \(f^2=1+\Delta i^2c^4\equiv1\pmod c\).
Equations \eqref{eq:cR}, \eqref{eq:cp}, and the normalized rank conclusion
verify every hypothesis of Lemma~\ref{lem:index}. We obtain
\begin{equation}
 p=R.                                       \label{eq:exactp}
\end{equation}
No sign of \(V\), and no restored value of \(h\), was used in this step.

\subsection{Ordered recovery of the index and scale}

Reducing \eqref{eq:q5carrier} modulo \(q^4\) gives \(q^4\mid r_0\).
Since \(P_0\equiv1\pmod{2XY^2}\), the first-Pell recurrence gives
\begin{equation}
 k\equiv2n_0\pmod{4XY^2}.                   \label{eq:firstmod}
\end{equation}
The modulus is a multiple of \(q^4\), because \(q^3\mid Y\).
Consequently
\[
 q^4\mid2n_0-R-1.
\]
But \(1\le n_0<p=R<q^4\), so
\(|2n_0-R-1|<q^4\). Hence
\begin{equation}
 2n_0=R+1.                                 \label{eq:exactn}
\end{equation}
Now \eqref{eq:firstmod} implies that
\[
 r_0=k-R-1=k-2n_0
\]
is a multiple of \(4XY^2\), hence of \(q^5\).
It is positive because \(n_0\ge2\) and
\(\psi_{P_0}(n_0)>n_0\).
Returning to the stronger congruence \eqref{eq:q5carrier}, we get
\(q^5\mid q^4X\), so \(q\mid X\).
The restored coordinates
\begin{equation}
 w=X/q>0,\qquad h=r_0/E>0                  \label{eq:restorewh}
\end{equation}
are integers; integrality of \(h\) follows from \(E=XY\mid4XY^2\).
This proves the removed index equation exactly.

\subsection{The main quotient and transport quotient are positive}

Put
\[
 \gamma_{\rm old}=\gamma+w-r_0/U.
\]
It is an integer, and \eqref{eq:Drecover} says
\[
 D=b_0+H\gamma_{\rm old}=X+ac+H\gamma_{\rm old}.
\]
Since \(D>(a+1)c\) and \eqref{eq:cR} gives \(c>kY>X\),
\[
 H\gamma_{\rm old}=D-ac-X>c-X>0.
\]
Thus \(\gamma_{\rm old}>0\).

We have already proved \(B_t>0\) and \((q-1)\mid B_t\).
Now \(X=qw\) gives
\[
 B_t=q\bigl((K_0+w)C-F\bigr).
\]
Since \(\gcd(q,q-1)=1\), there is a positive integer
\[
 z_{\rm new}=\frac{(K_0+w)C-F}{q-1}.
\]
Set \(z=z_{\rm new}+1>0\). Then
\[
 (K_0+w)C+(q-F)-z(q-1)=1,
\]
the original transport condition. Finally set \(T_{\rm old}=fT>0\).
Its source auxiliary expression is exactly
\[
 c(T_{\rm old}f-1)-Rf^2=f^2(cT-R)-c=V.
\]
Together with \(s=Y/q^3>0\), \(\tau_{\rm old}=|\tau|>0\), and
\(y=|V+u|>0\), we have reconstructed every positive source coordinate.
Theorem~\ref{thm:source} gives a finite ZFC refutation.

\section{Completeness and degree of the thirteen-variable equation}
\label{sec:complete13}

\subsection{Constructing all new positive coordinates}

Start from a finite refutation and choose the canonical source witnesses
of Section~\ref{sec:canonical}. Keep the normalized \(i,f\) and the
positive outer coordinates. Since \(B\equiv-1\pmod3\) and \(N\) is odd,
\[
 J=1+B+\cdots+B^{N-1}\equiv1\pmod3.
\]
As \(J>1\), the new \(j=(J-1)/3\) is a positive integer.

The canonical \(X=2^R\) and \(R>3q\), \(q\ge16\), give
\(w=X/q>q\). From the source transport equation,
\[
 z(q-1)=(K_0+w)C+q-F-1>w+1>q,
\]
where \(C\ge1\) and \(q-F=C+\alpha\ge2\).
Thus \(z\ge2\). Define
\begin{equation}
 t=q(z-1)\ge q,\qquad
 r=Wt+\tau-\ell-k-s,\quad s=Y/q^3.          \label{eq:rcomplete}
\end{equation}
The old transport identity gives \(B_t=(q-1)t\).
Since \(\tau>\ell\) and \(k>XY>s\),
\[
 r>qW-k-s>qW-2k
      =\bigl(q(q-1)-2\bigr)k+q>0.
\]
All terms in \eqref{eq:rcomplete} are integers. Substitution gives
\(Z=M_0\tau\), so \(A_0=0\).

The source first index is \(n_0=(R+1)/2\), so
\eqref{eq:firstmod} makes \(r_0=k-R-1\) a positive multiple of
\(4XY^2\). For the canonical scales this is a multiple of \(q^7\),
and therefore of \(q^5\). Set
\begin{equation}
 \gamma=\gamma_{\rm old}-w+r_0/q^5.         \label{eq:gammacomplete}
\end{equation}
This is a positive integer by \eqref{eq:canonicalgamma}.
Its substitution yields \(D_*=q^5D\), hence \(A_1=0\).

Finally apply Lemma~\ref{lem:auxcomplete} with \(S=\Delta I\).
It supplies positive \(T,u\), with \(V=V_*\) and \(V+u=y_*\).
This changes no carrier coordinate. The strong and auxiliary residuals
vanish, and all thirteen supplied coordinates are positive.
This proves the reverse existence implication. Combining it with
Section~\ref{sec:sound13} and negating establishes the first statement of
Theorem~\ref{thm:main}.

\subsection{Exact cancellation and degree}
\label{sec:degree13}

Let \(q_0=3\beta j\), \(k_0=\eta+\zeta\), and define the homogeneous forms
\[
 R_0=q_0^3(q_0-F),\qquad
 B_{t,0}=(K_0q_0+X)(q_0-F-\alpha)-q_0F.
\]
For the first carrier, write
\[
 T_0=Z-M_0\ell=M_0(k+r)+(q-1)Y-q^3WB_t.
\]
The exact identity
\begin{equation}
 A_0=2M_0\ell T_0+T_0^2-M_0^2\ell k-M_0^2
                                                        \label{eq:A0cancel}
\end{equation}
removes the apparent degree-sixteen terms. Its degree-fifteen part is
\begin{equation}
 \ini_{15}(A_0)=-2q_0^8k_0^2XY^2B_{t,0}.    \label{eq:headA0}
\end{equation}
For the main carrier put \(T_1=U\gamma+q^4X-r_0\).
The exact identity
\begin{equation}
 \begin{split}
 A_1={}&U^2(2acX+X^2-Hc^2-1)\\
       &+2U(ac+X)HT_1+H^2T_1^2
 \end{split}                               \label{eq:A1cancel}
\end{equation}
gives degree \(17\), with
\begin{equation}
 \ini_{17}(A_1)=8q_0^{10}k_0X^2Y^3\gamma.   \label{eq:headA1}
\end{equation}
The other residuals have degrees \(13\) and \(14\):
\[
 \ini_{13}(A_a)=2X^2Y^2f^4uR_0,\qquad
 \ini_{14}(A_s)=-i^2k_0^4X^2Y^6.
\]
Thus only \(A_1^2\) reaches degree \(34\), and
\begin{equation}
 \ini_{34}(P_{13})
   =64(3\beta j)^{20}(\eta+\zeta)^2X^4Y^6\gamma^2.
                                                        \label{eq:head13}
\end{equation}
The coefficient of
\(j^{20}X^4Y^6\eta^2\gamma^2\) is \(64(3\beta)^{20}>0\).
This proves exact degree after specialization to the fixed coefficient
tuple, not just a generic formal-port upper bound.

\subsection{The literal evaluator}

The supplied evaluator has \(52\) multiplications, \(24\) additions, and
\(17\) subtractions, for \(93\) live binary operations.
One saving is the exact factoring identity
\[
 T_{\rm tr}=(qK_0+X)C+q(q-F-1),\qquad
 T_{\rm tr}-q(q-1)=B_t.
\]
Consequently the less economical numerator
\[
 M_0(\ell+k+qW+r)+(q-1)Y-q^3WT_{\rm tr}
\]
equals \(Z\) identically. The independent checker retains this unfactored
expression when comparing to the emitted graph, so it checks the rewrite
rather than repeating it.

Two further shared-register rewrites preserve the exact polynomial.
The register \(f^2-1\), already needed by \(K_a\), is reused in
\(A_s=(f^2-1)-\Delta I^2\). For the auxiliary residual, the evaluator
puts \(y_0=V+u\) and computes
\[
 A_a=y_0^2-K_a(y_0^2-V^2)-1.
\]
The identity \(y_0^2-V^2=u(2V+u)\) shows that this is exactly
\eqref{eq:Aa}. These are arithmetic-cost changes; the four residuals
and their degree certificates are unchanged.

\section{A divided norm gives twelve positive unknowns}
\label{sec:p12}

Supply exactly
\begin{equation}
 j,F,\alpha,i,T,X,Y,r,\eta,\zeta,n,\gamma\in\Zpos.
                                                        \label{eq:vars12}
\end{equation}
Use the common registers \eqref{eq:common13}--\eqref{eq:firstcarrier},
and \(A_0,A_1\) from \eqref{eq:A0}--\eqref{eq:A1}, with the single
change
\begin{equation}
 J=6j+1.                                    \label{eq:J12}
\end{equation}
The expressions \(K_a,V,A_a,A_s\) of the thirteen-variable construction
are replaced by the following explicit polynomial registers:
\begin{align}
 D_f&=1+\Delta I^2,&
 C_*&=\Delta ic,& M&=\Delta ic^3, \nonumber\\
 Q&=MC_*=(\Delta I)^2,& d_*&=Q-1,&
 V&=D_f(cT-R)-c,                             \label{eq:norm12reg}\\
 A_*&=d_*M(n^2+D_f)-2d_*n+C_*(1-V^2),        \label{eq:Anorm}\\
 B_*&=2d_*(Mn-1),&
 N_*&=A_*^2-B_*^2D_f.                       \label{eq:Nnorm}
\end{align}
The twelve-variable polynomial is
\begin{equation}
 \boxed{P_{12}=(1+A_0)(1+A_1)(1+N_*)-1.}     \label{eq:P12}
\end{equation}
It has \(104\) live binary operations. All powers in these definitions
are fixed, and no division occurs in the polynomial.

\subsection{The exact identity behind the divided norm}

For a proof-only \(f\) satisfying \(f^2=D_f\), put \(t=n-f\).
Direct expansion gives
\begin{align}
 A_*-B_*f
   &=d_*Mt^2-2d_*t+C_*(1-V^2),              \label{eq:dividedfirst}\\
 M(A_*-B_*f)
   &=d_*\bigl(M(n-f)-1\bigr)^2-QV^2+1.      \label{eq:dividedidentity}
\end{align}
For \eqref{eq:dividedidentity}, multiply \eqref{eq:dividedfirst} by
\(M\), use \(MC_*=Q\), and \(d_*=Q-1\).
The relation is an exact polynomial identity after substituting
\(D_f=f^2\); it is not a numerical approximation or a formal assertion
about independent square classes.

The factorization \(Q=MC_*\), with
\(M=\Delta ic^3\), \(C_*=\Delta ic\), distributes the degree between the
two terms of \(A_*\). The residue of the canonical auxiliary ordinate
permits division by \(M\) in the completeness proof. In the actual
polynomial, that quotient is encoded by the supplied coordinate \(n\)
and the cleared identity above.

\subsection{The product finalizer recovers the three equations}

The integer \(\beta\) is odd, and \(J=6j+1\) is odd, so \(q\) is even.
Thus \(M_0=q^3(q-1)\) and \(U=q^5\) are even. It follows that
\[
 A_0\equiv Z^2\pmod4,\qquad
 A_1\equiv D_*^2\pmod4.
\]
Also \(B_*\) is even, hence \(N_*\equiv A_*^2\pmod4\).
At a zero of \eqref{eq:P12}, each of its three integer factors is
\(+1\) or \(-1\). A factor \(-1\) would make its residual
\(-2\equiv2\pmod4\), impossible for a square modulo four.
Therefore
\begin{equation}
 P_{12}=0
 \quad\Longleftrightarrow\quad A_0=A_1=N_*=0
 \qquad\hbox{on }\Zpos^{12}.                 \label{eq:unitfinal}
\end{equation}
No nonnegativity of the individual residuals is asserted.

\subsection{Every twelve-variable zero has a positive thirteen-variable lift}

For every positive input, \(c\ge3\), \(\Delta>0\), \(Q>1\), \(M>1\),
and \(D_f>c+1\). In particular
\[
 B_*=2(Q-1)(Mn-1)>0.
\]
There is no zero-denominator branch. If \(N_*=0\), the rational number
\[
 f=A_*/B_*
\]
has square \(D_f\); by Lemma~\ref{lem:rational} it is a nonzero integer.
Set
\[
 y=M(n-f)-1.
\]
Identity \eqref{eq:dividedidentity} gives
\begin{equation}
 QV^2-(Q-1)y^2=1.                           \label{eq:recover12aux}
\end{equation}
Both \(f\) and \(y\) may have either sign, but a positive lift exists
without invoking a new signed source theorem.

Indeed, let \(b=cT-R\in\Z\). The unconditional identity \(V=D_fb-c\)
and \(D_f>c+1\) imply
\[
 b\le0\ \Longrightarrow\ V\le-c<-1,\qquad
 b\ge1\ \Longrightarrow\ V\ge D_f-c>1.
\]
Thus \(|V|>1\) on every positive input. Equation
\eqref{eq:recover12aux} now gives
\[
 (Q-1)(y^2-V^2)=V^2-1>0,
\]
so \(|y|>|V|\). Define the three coordinates of a thirteen-variable lift
by
\begin{equation}
 j_{13}=2j,\qquad f_{13}=|f|,\qquad u_{13}=|y|-V.
                                                        \label{eq:lift12to13}
\end{equation}
All are strictly positive. The identity
\(3j_{13}+1=6j+1\) makes every outer register identical.
Since \(f_{13}^2=D_f\),
\[
 f_{13}^2-\Delta I^2=1,\qquad
 \Delta(f_{13}^2-1)=Q,\qquad
 f_{13}^2(cT-R)-c=V.
\]
Also \(V+u_{13}=|y|\), so its auxiliary residual vanishes by
\eqref{eq:recover12aux}. The two carriers already vanish by
\eqref{eq:unitfinal}. We have constructed an actual positive zero of
\(P_{13}\), whose soundness was proved above. This proves soundness of
\(P_{12}\), including all signed-root branches.

\subsection{Canonical completeness and the ordinate residue}

Start with canonical source witnesses. The repunit \(J\) is odd because
\(B\) is even, and \(J\equiv1\pmod3\) by the odd canonical exponent.
Thus \(J\equiv1\pmod6\), and \(j=(J-1)/6>0\) is an integer.
Choose \(r,\gamma\) exactly as in \eqref{eq:rcomplete} and
\eqref{eq:gammacomplete}; these give \(A_0=A_1=0\).

Use Lemma~\ref{lem:auxcomplete} to reselect the canonical auxiliary
pair \(V_*,y_*\) and quotient \(T\). It gives
\[
 y_*=\psi_{\Delta I}(R)\equiv-1\pmod{(\Delta I)^2}
                                  =-1\pmod Q.
\]
Since \(M\mid Q\), the number
\[
 t=(y_*+1)/M
\]
is a positive integer. Take
\begin{equation}
 n=f+t>0.                                  \label{eq:ncomplete12}
\end{equation}
Now \(y_*=M(n-f)-1\), so the auxiliary Pell identity and
\eqref{eq:dividedidentity} give \(M(A_*-B_*f)=0\).
Because \(M>0\), \(A_*=B_*f\), whence \(N_*=0\).
All twelve supplied coordinates are positive, and
\eqref{eq:P12} vanishes.

The residue assertion is an all-index consequence of
Lemma~\ref{lem:odd}; finite Pell examples are not used as its proof.
It is needed only for the freshly chosen canonical auxiliary ordinate.
Together with the positive lift \eqref{eq:lift12to13}, this completes
the second equivalence in Theorem~\ref{thm:main}.

\section{Exact degree and cost of the twelve-variable equation}
\label{sec:degree12}

For this section put
\[
 q_0=6\beta j,\quad k_0=\eta+\zeta,\quad c_0=k_0Y,\quad
 \Delta_0=X^2Y^2,\quad R_0=q_0^3(q_0-F),
\]
and
\[
 B_{t,0}=(K_0q_0+X)(q_0-F-\alpha)-q_0F,\qquad
 D_{f,0}=\Delta_0 i^2c_0^4.
\]
The carrier computations \eqref{eq:headA0}--\eqref{eq:headA1} still
apply with this \(q_0\). In particular their exact degrees are \(15,17\).
The relevant degrees for the norm are
\[
 \deg D_f=14,\quad \deg C_*=7,\quad \deg M=11,\quad
 \deg Q=\deg d_*=18,\quad \deg V=18.
\]
The two degree-\(43\) contributions to \(A_*\) are \(d_*MD_f\)
and \(-C_*V^2\). All its other summands have degree at most \(32\).
Their combined leading form is
\begin{align}
 A_{*,43}
  &=C_{*,0}D_{f,0}^{\,2}(\Delta_0c_0^2-R_0^2) \nonumber\\
  &=i^5k_0^9X^6Y^{15}
               \bigl(k_0^2X^2Y^4-R_0^2\bigr).             \label{eq:headnormA}
\end{align}
This is nonzero. Meanwhile \(\deg B_*=30\), so
\(\deg(B_*^2D_f)=74<86\). Hence
\[
 \deg N_*=86,\qquad \ini_{86}(N_*)=A_{*,43}^{\,2}.
\]
The product finalizer therefore has exact degree \(15+17+86=118\),
with leading form
\begin{equation}
 \boxed{\begin{aligned}
 \ini_{118}(P_{12})={}&-16q_0^{18}i^{10}k_0^{21}
           X^{15}Y^{35}\gamma B_{t,0}\\[-2pt]
 &{}\times\bigl(k_0^2X^2Y^4-R_0^2\bigr)^2.
 \end{aligned}}                                             \label{eq:head12}
\end{equation}
For a coefficient independent of the other three fixed ports, take
\begin{equation}
 j^{18}i^{10}\eta^{25}X^{20}Y^{43}\gamma\alpha.
                                                        \label{eq:monomial12}
\end{equation}
Its coefficient is exactly
\begin{equation}
 16(6\beta)^{18}>0.
                                                        \label{eq:coefficient12}
\end{equation}
Indeed, the minimum \(j\)-power \(18\) forces the choice of
\(-X\alpha\) from \(B_{t,0}\) and of
\(k_0^4X^4Y^8\) from the squared bracket.
No specialization of \(K_0,M_C,M_F\) can cancel this coefficient.

The literal graph has \(60\) multiplications, \(26\) additions, and
\(18\) subtractions, totaling \(104\). It shares \(ic=i\,c\), using
\[
 C_*=\Delta(ic),\qquad
 D_f=1+(\Delta c^2)(ic)^2.
\]
The register \(\Delta c^2\) is already used by \(A_1\). These identities
save a multiplication without changing either polynomial.
For an alternative finalizer, the same three residuals give
\begin{equation}
 P_{12}^{\rm sos}=A_0^2+A_1^2+N_*^2.          \label{eq:P12sos}
\end{equation}
It has the same positive zeros, exact degree \(172\), and \(103\)
binary operations. Its highest homogeneous form is \(A_{*,43}^{\,4}\).
The one-operation saving therefore comes with a larger degree.
The archive retains this companion as
\sourcepath{positive/p12\_sos\_split2.json}.

\section{Two equations over unrestricted integers}
\label{sec:integer}

The integer constructions take a different route. They begin with the
literal forty-five-integer quartic system \(Q_{45}\) of source 57
\cite{quartic45} and eliminate variables that occur in monic defining
equations. They retain the source's entire
\emph{mixed-parity-unrounded-1064921} tuple, with
\(B=2^{1064921}\). This is different from the compact tuple used by
\(P_{13},P_{12}\); neither family borrows individual ports from the other.

\subsection{The full source system}

All variables in this section are unrestricted integers. Define
\begin{equation}
 \pos(x_1,x_2,x_3)=1+x_1^2+x_2^2+x_3(x_3+1).
                                                        \label{eq:Pos}
\end{equation}
This is always strictly positive and represents every positive integer.
For the latter assertion, apply Legendre's three-square theorem to
\(4n-3\), which is \(1\) modulo four and is not of the excluded form
\(4^a(8b+7)\). One square is odd and the other two are even;
writing them as \((2x_3+1)^2,(2x_1)^2,(2x_2)^2\) gives
\eqref{eq:Pos}.

The source's thirty-two arithmetic coordinates are
\[
 \begin{gathered}
 J,F,\alpha,z,f,h,r,v,w,\tau,\eta,\zeta,y,\gamma,o,j,\\
 N,E,P,c,a,X,D,C_2,L,F_2,V,g,I,B_a,\Delta,u.
 \end{gathered}
\]
They are accompanied by the thirteen guard coordinates consisting of
four triples
\(\boldsymbol e,\boldsymbol z,\boldsymbol b,\boldsymbol d\in\Z^3\)
and one integer \(v_q\). Set
\[
 q=\beta J+1,\quad k=\eta+\zeta,\quad H=4a+3,\quad C=q-F-\alpha,
\]
and the quadratic expression
\[
 \mathcal R=qN-qL-q^2+F(q+L-N)+\alpha L+(M_C+qM_F)J.
\]
The twenty-eight source residuals are the left sides of the following
equations:
\begin{align}
 L-(q^2-1)&=0,& N-q(L+1)&=0, \nonumber\\
 X-wq&=0,& a-E-(2v+1)N&=0, \nonumber\\
 E-X(a-E)&=0,& c-\eta-k(a-E)&=0, \nonumber\\
 P-E(c-\eta)-k&=0,& \tau^2-P(P-k)-1&=0, \nonumber\\
 D-X-ac-\gamma H&=0,& \Delta-a^2-H&=0, \nonumber\\
 C_2-c^2&=0,& D^2-\Delta C_2-1&=0, \nonumber\\
 k-hE-\mathcal R-1&=0,& (K_0+w)C-F-z(q-1)&=0, \nonumber\\
 F_2-f^2&=0,& V+c-fo&=0, \nonumber\\
 V-cj+k-hE-1&=0,& V^2+g-y^2&=0, \nonumber\\
 I-(2r+1)C_2&=0,& B_a-I^2&=0, \nonumber\\
 B_a-\Delta(F_2-1)&=0,& y^2-B_ag-1&=0, \nonumber\\
 w-u^2&=0,& \eta-\pos(\boldsymbol e)&=0, \nonumber\\
 \zeta-\pos(\boldsymbol z)&=0,& q-2v_q^2&=0, \nonumber\\
 F\alpha-\pos(\boldsymbol b)&=0,&
 FC-\pos(\boldsymbol d)&=0.                 \label{eq:Q45system}
\end{align}
These are the literal source residuals, with renamed guard coordinates
and harmless regrouping of terms. In particular, the transport
coordinate here is the shifted quotient satisfying
\((K_0+w)C=F+z(q-1)\).
Under the first two equations,
\(\mathcal R=((q-1)(q-F)+\alpha)(q^2-1)+(M_C+qM_F)J\).
Every displayed residual has degree at most two.

The source theorem states that the sum of their squares has no integer
zero if and only if \(\Con(\ZFC)\), provably in ZFC
\cite{quartic45,oddscale}. Its domain proof uses the dyadic guard
\(q=2v_q^2\) to force \(J>0\); the two product guards force
\(F,\alpha,C>0\), since these three nonzero integers have the same sign
and sum to \(q>0\). The inherited odd-scale and signed-root arguments
then restore the unguarded source coordinates. In particular they give
\(X=2^R\) and \(w>F\), so the literal transport residual above implies
\(z>0\). Canonical completeness supplies the dyadic guard from the source
tuple's odd binary exponent. These unbounded source arguments remain
inherited; the reductions below alter none of them.

\subsection{The exact elimination principle}

\begin{lemma}[Unit elimination]
\label{lem:eliminate}
Suppose an integer polynomial system contains
\[
 x-A(\boldsymbol y)=0
\]
with \(A\in\Z[\boldsymbol y]\), and other residuals
\(F_i(x,\boldsymbol y)\). Deleting that defining residual and
substituting \(x=A(\boldsymbol y)\) in every other residual gives a
system whose integer solutions are in bijection with the original ones.
\end{lemma}
\begin{proof}
The maps are projection \((x,\boldsymbol y)\mapsto\boldsymbol y\)
and reconstruction
\(\boldsymbol y\mapsto(A(\boldsymbol y),\boldsymbol y)\).
They are inverse on the two solution sets by direct substitution.
No division, root extraction, sign premise, or nonvanishing condition
is involved. The same assertion holds when the defining residual has
the opposite sign.
\end{proof}

Iterating the lemma along a triangular list preserves the exact integer
solution set up to a specified bijection. Replacing the remaining system
by the sum of its squared residuals preserves its zeros over \(\Z\).

\subsection{The thirty-seven-variable sextic}

In \eqref{eq:Q45system}, eliminate the following eight coordinates using
their displayed source definitions, substituting everywhere:
\begin{equation}
 \begin{aligned}
 B_a&=I^2,& C_2&=c^2,& F_2&=f^2,& N&=q(L+1),\\
 w&=u^2,& X&=D-ac-\gamma H,&
 \eta&=\pos(\boldsymbol e),&
 \zeta&=\pos(\boldsymbol z).
 \end{aligned}                              \label{eq:elim37}
\end{equation}
The remaining \(37\) source variables and \(20\) residuals define
\(Q_{37}\) as their sum of squares. Equations
\eqref{eq:Q45system} and \eqref{eq:elim37} are a complete polynomial
specification; \sourcepath{integer/sextic37.json} lists the resulting
residuals without requiring the reader to perform the substitutions.

Every residual now has degree at most three. The useful choice is the
definition of \(X\): the main-root equation expresses it quadratically,
whereas combining \(X=wq\) with \(w=u^2\) would make it cubic.
The other equation for \(X\) becomes
\[
 D-ac-\gamma H=u^2q,
\]
which is cubic, and the substituted relation \(E=X(a-E)\) remains cubic.
The substituted positivity polynomials are quadratic and enter the
relevant source residuals with at most one additional linear factor.

For an exact lower bound, the residual
\[
 I-(2r+1)c^2
\]
contains the only occurrence of \(r\) in the new system. Its square
contributes coefficient \(4\) to \(r^2c^4\), and no other square
contributes to that monomial. Therefore
\[
 \deg Q_{37}=6
\]
on every fixed coefficient slice, including the authentic source tuple.
The full symbolic expansion has \(718\) monomials when the four ports
are retained as formal symbols. Lemma~\ref{lem:eliminate} gives the
exact solution bijection with \(Q_{45}\), proving the consistency
equivalence.

\subsection{The thirty-four-variable octic}

For the second elimination, use the following triangular substitutions:
\begin{equation}
 \begin{aligned}
 L&=q^2-1,& N&=q^3,& X&=wq,\\
 D&=wq+ac+\gamma H,& \Delta&=a^2+H,& C_2&=c^2,\\
 F_2&=f^2,& B_a&=I^2,& g&=y^2-V^2,\\
 \eta&=\pos(\boldsymbol e),&
 \zeta&=\pos(\boldsymbol z).
 \end{aligned}                              \label{eq:elim34}
\end{equation}
There are \(34\) remaining integer variables and \(17\) residuals,
all of degree at most four. Define \(Q_{34}\) as their sum of squares.
The explicit residual list is \sourcepath{integer/octic34.json}.
For example, two surviving equations are
\[
 y^2-I^2(y^2-V^2)-1=0,\qquad
 I^2-(a^2+4a+3)(f^2-1)=0.
\]
The second has coefficient \(-1\) on \(a^2f^2\); the other residuals
have smaller degree in the pair \(a,f\). Thus the coefficient of
\(a^4f^4\) in their sum of squares is exactly \(1\).
It follows that \(\deg Q_{34}=8\), independently of the fixed ports.
Its symbolic sum of squares has \(814\) monomials.

Again Lemma~\ref{lem:eliminate} gives an exact integer-solution bijection
with \(Q_{45}\). Together with the inherited source theorem, this proves
the two integer assertions in Theorem~\ref{thm:main}. The source guards
remain algebraically present after substitution; no positive coordinate
has been discarded on the strength of its name.

\subsection{Comparison and scope of the integer reductions}

The historical sextic and octic retained in the coefficient dossier
have respectively \(42\) variables and \(21\) residuals, and \(36\)
variables and \(15\) residuals. Their exact degrees are \(6\) and \(8\).
Thus the new sextic saves five integer variables, and the octic saves
two. The latter has two more residual equations. These comparisons are
between independently calibrated equations with matching domains and
degrees; no direct solution bijection with those historical comparison
polynomials is claimed.

The actual new bijections are to the fixed \(Q_{45}\) source. Its literal
packet has SHA-256
\begin{center}
\scriptsize\texttt{2b4d3d372fd899042eb3db9ee25e085193204640d28009807fde0778bc9528ac}.
\end{center}
The source and its coefficient dossier are included, so both the precise
polynomial and the binding of its four fixed ports can be checked.

\section{Fixed coefficients and reproducible verification}
\label{sec:binding}

\subsection{The complete compact binding for the positive equations}

The main binding is \sourcepath{positive/ACTIVE\_BINDING.json}; its source
is the complete source-44 manifest
\sourcepath{coefficients/compact44/COEFFICIENT\_BINDING.json}.
The numerical ports are the unsigned little-endian integers in the four
decompressed files below. Explicitly, bytes \(b_0,\ldots,b_{t-1}\)
represent \(\sum_{\nu=0}^{t-1}b_\nu256^\nu\).

\begin{center}
\begin{tabular}{@{}llrrr@{}}
\toprule
Port & File in \sourcepath{compact44/coefficients/}
 & Bits & Raw bytes & Gzip bytes\\
\midrule
\(\beta\) & \sourcepath{beta.bin.gz} & 590069 & 73759 & 108\\
\(K_0\) & \sourcepath{K.bin.gz} & 592884 & 74111 & 8940\\
\(M_C\) & \sourcepath{MC.bin.gz} & 590069 & 73759 & 193\\
\(M_F\) & \sourcepath{MF\_source.bin.gz} & 590070 & 73759 & 204\\
\midrule
Total & & & 295388 & 9445\\
\bottomrule
\end{tabular}
\end{center}

The raw SHA-256 digests, in the same order, are:
\begin{quote}
\scriptsize\ttfamily
0664a868cbfef9869ea05c40a1c445563f685ff857092594c1c1223bc4b21f9d\\
c6d523bd32af8dd0345db843cdadf6932bfcad186912a6d112314b60657c554e\\
6007cabd5f5228618e8fb613f386d8fc3692a24520d6e4cd72af0262aae1ade4\\
48333b0d30a4edeb8579f44286cb59806c5622cb42b367f5b5f9ee827711c737
\end{quote}
The files and deterministic builder, rather than truncated printed
digits, specify the complete fixed integers. These digests authenticate
that specification. In particular \(\beta=2^{590069}-1\).
The four ports are bound together and the \(M_F\) payload already contains
the source shift.

The retained builder \sourcepath{build\_coefficients.py}, ordered
\sourcepath{k215\_transitions.json}, and adapter
\sourcepath{PROOF.md} are in \sourcepath{coefficients/compact44/}.
The compact native encoding uses radix \(2^{67}\), layout length \(8807\),
and \(d=67\cdot8807=590069\), with padding prime \(581197\)
\cite{source44}. Its finite coefficient audits check the implemented
masks and carry/controller arithmetic. Their finite scope does not
replace the unbounded native-language and calibration theorems.

\subsection{The separate integer binding}

For \(Q_{37},Q_{34}\), the authoritative file is
\sourcepath{integer/coefficients\_source/ACTIVE\_BINDING.json}.
Its ports have the following bit lengths:
\[
 \begin{array}{c|rrrr}
 \text{port}&\beta&K_0&M_C&M_F\\ \hline
 \text{bits}&1064921&1065043&1064921&1064922.
 \end{array}
\]
Here \(\beta=2^{1064921}-1\), with the already shifted high-mask port.
The complete raw payloads occupy \(532479\) bytes; their supplied gzip
forms occupy \(10466\) bytes. The dossier retains the package-34
coefficient builder, independent decoder, ordered Wang943 input,
literal machine input, source proof, historical comparator packets,
and source-transfer ledger.

The recorded reconstruction reproduces the raw payloads, gzip payloads,
and full historical binding manifest exactly. This establishes a
complete numerical binding for these equations; it does not transfer the
compact source-44 tuple to an integer source whose domain theorem was
proved for a different tuple.

\subsection{What the exact checks establish}

The main archive entry point is
\begin{center}
\sourcepath{python3 verify\_all.py}.
\end{center}
The individual builders and checks are retained under
\sourcepath{positive/} and \sourcepath{integer/}.
For the two positive-variable equations, the checks read literal graphs
as data and compare them with independently specified formulas.

For \(P_{13}\), the checker expands the four residuals exactly over
the thirteen unknowns and four symbolic fixed ports. It verifies
\eqref{eq:A0cancel}, \eqref{eq:A1cancel}, the exact residual degrees,
the leading form \eqref{eq:head13}, all live inputs, and the operation
ledger. The supplied circuit is thereby tied to the mathematical
specification, including its nonzero coefficient after actual-port
specialization.

For \(P_{12}\), the checker expands the prefix containing the two
carriers and the norm's input registers. It checks the norm suffix
as an exact identity over independent cut registers
\(D_f,M,C_*,V,n\), and the finalizer over independent
\(A_0,A_1,N_*\). These identities compose to an exact check of the
whole graph. It separately verifies \eqref{eq:dividedfirst} and
\eqref{eq:dividedidentity}, the leading form \eqref{eq:head12}, and
the selected coefficient \eqref{eq:coefficient12}.
No random evaluation is substituted for an identity.

For the two integer equations, the builder checks that each eliminated
residual is literally plus or minus the claimed unit definition.
It performs exact sparse substitution, expands the full sum of squares,
and verifies the variable counts, degrees, liveness, and coefficients
\([r^2c^4]Q_{37}=4\) and \([a^4f^4]Q_{34}=1\).
The new checks use explicit exceptions rather than removable
assertions and can also be run under \sourcepath{python3 -O}.
Historical programs that require a normal Python invocation retain
that requirement in their own documentation.

\subsection{Mathematical and formal boundaries}

There are three distinct kinds of evidence in the archive. Exact finite
algebra verifies that the delivered formulas, literal evaluators,
substitutions, degrees, and numerical ports agree. The mathematical
proofs in this paper establish the new unbounded existence reductions,
including every recovery and completeness branch. The inherited native
compiler and syntactic calibration identify the old source equations
with finite ZFC refutations.

The theorem follows by composing these mathematical implications and
negating the resulting existential statements. No proof of
\(\Con(\ZFC)\), no example of a refutation, and no assumption deciding
whether one exists enters the construction of a coefficient.
The ordinary proofs are formalizable in ZFC, but this archive does not
claim a newly compiled proof-assistant endpoint for the entire chain.
Internal separately conducted reviews are research checks within this
collaboration, not external human endorsements.

The inherited source material and its original notices are retained.
The source licenses include ProveIt's MIT No Attribution terms and the
retained Wade license; the native compiler and machine lineage remain
attributed within their original scopes. The changes proved here are
the combined congruence carriers, their recovery order and positive
completeness construction, the divided norm specialization, and the two
exact integer eliminations. No new universal variable bound or global
minimum is asserted.

\begin{thebibliography}{99}
\small
\bibitem{source44}
ProveIt contributors,
\emph{Complement-packed native clauses: a smaller complete coefficient tuple},
source 44, native proof and complete numerical binding.
\href{https://github.com/VladimirReshetnikov/ProveIt/blob/beb0192cb8dc36f78c369c080f2adfbc1887a486/Computability/DiophantineConsistency/con-zfc-diophantine-forms/44-b4-sparse/native/PROOF.md}
{\texttt{44-b4-sparse/native/PROOF.md}}.

\bibitem{ordinary}
ProveIt contributors,
\emph{Ordinary strong norms and signed witnesses for the blank computation},
especially Sections 4--7.
\href{https://github.com/VladimirReshetnikov/ProveIt/blob/beb0192cb8dc36f78c369c080f2adfbc1887a486/Computability/DiophantineConsistency/con-zfc-diophantine-forms/35-cancel12/sources/catalogue/25-hybrid12/upstream/ORDINARY_BLANK_PROOF.md}
{\texttt{35-cancel12}, preserved \texttt{ORDINARY\_BLANK\_PROOF.md}}.

\bibitem{pellaudit}
ProveIt contributors,
\emph{Pell dependency audit for the 67-operation blank-input polynomial},
especially Sections 3.1 and 4.
\href{https://github.com/VladimirReshetnikov/ProveIt/blob/beb0192cb8dc36f78c369c080f2adfbc1887a486/Computability/DiophantineConsistency/con-zfc-diophantine-forms/36-small-tuple/proofs/INHERITED_PELL_AUDIT.md}
{\texttt{36-small-tuple/proofs/INHERITED\_PELL\_AUDIT.md}}.

\bibitem{source38}
ProveIt contributors,
\emph{Eleven positive variables: normalized Pell recovery},
especially Section 5.
\href{https://github.com/VladimirReshetnikov/ProveIt/blob/beb0192cb8dc36f78c369c080f2adfbc1887a486/Computability/DiophantineConsistency/con-zfc-diophantine-forms/38-batching11/eleven/ELEVEN_POSITIVE_PROOF.md}
{\texttt{38-batching11/eleven/ELEVEN\_POSITIVE\_PROOF.md}}.

\bibitem{source51}
ProveIt contributors,
\emph{Squared-modulus auxiliary reconstruction}.
\href{https://github.com/VladimirReshetnikov/ProveIt/blob/beb0192cb8dc36f78c369c080f2adfbc1887a486/Computability/DiophantineConsistency/con-zfc-diophantine-forms/51-sparse-arith/short/SQUARED_MODULUS_PROOF.md}
{\texttt{51-sparse-arith/short/SQUARED\_MODULUS\_PROOF.md}}.

\bibitem{source57}
ProveIt contributors,
\emph{Two separate quadratic norms: twelve positive variables at degree 242},
especially Sections 5 and 9.
\href{https://github.com/VladimirReshetnikov/ProveIt/blob/beb0192cb8dc36f78c369c080f2adfbc1887a486/Computability/DiophantineConsistency/con-zfc-diophantine-forms/57-separated-degrees/low_degree/PROOF.md}
{\texttt{57-separated-degrees/low\_degree/PROOF.md}}.

\bibitem{source58}
ProveIt contributors,
\emph{Integral-root carriers}, report, especially the inherited kernel
and canonical completeness in Sections 4 and 6.
\href{https://github.com/VladimirReshetnikov/ProveIt/blob/beb0192cb8dc36f78c369c080f2adfbc1887a486/Computability/DiophantineConsistency/con-zfc-diophantine-forms/58-integral-root/report/report.tex}
{\texttt{58-integral-root/report/report.tex}}.

\bibitem{source59}
ProveIt contributors,
\emph{Thirteen-positive degree-separated adjunct}.
\href{https://github.com/VladimirReshetnikov/ProveIt/blob/beb0192cb8dc36f78c369c080f2adfbc1887a486/Computability/DiophantineConsistency/con-zfc-diophantine-forms/59-degree-separated/positive13_adjunct/PROOF.md}
{\texttt{59-degree-separated/positive13\_adjunct/PROOF.md}}.

\bibitem{quartic45}
ProveIt contributors,
\emph{A quartic consistency equation in 45 unrestricted integer variables}.
\href{https://github.com/VladimirReshetnikov/ProveIt/blob/beb0192cb8dc36f78c369c080f2adfbc1887a486/Computability/DiophantineConsistency/con-zfc-diophantine-forms/57-separated-degrees/quartic/QUARTIC45_PROOF.md}
{\texttt{57-separated-degrees/quartic/QUARTIC45\_PROOF.md}}.

\bibitem{oddscale}
ProveIt contributors,
\emph{Odd-scale integer lifts and their domain reduction}.
The historical proof and its exact source identities are retained in
\texttt{integer/coefficients\_source/low\_degree/upstream/}
\texttt{ODD\_SCALE\_LIFTS\_PROOF.md}, with the coefficient dossier's
transfer ledger and preserved comparison packets.
\end{thebibliography}

\bigskip
\footnotesize
All linked catalogue references above are pinned to the comparison
revision
\begin{center}\snapshot.\end{center}
Earlier commits appearing inside retained source
documents describe their own historical dependency captures; they do
not change the comparison pin or the active numerical bindings of this
archive.
\end{document}
